\section{Materials and Methods}
\label{sec:methods}

\subsection{Scenario construction and simulator credibility}

The public pipeline generated correction scenarios at three frozen levels of
physical detail. F0 uses central Earth point-mass gravity. F1 adds the Moon and
Sun as third-body gravitational sources. F2 additionally includes Earth J2,
the leading correction for Earth's equatorial bulge, and tighter numerical
integration settings. Across these levels the implementation retained the
configured maneuver model, mass depletion, event handling, crew-related
constraints, uncertainty variables, and numerical feasibility checks. The
intent was to predict benchmark-defined cost or feasibility outcomes, not to
control a spacecraft. Source files and scenario construction were frozen before
the model campaigns; later reports could summarize outputs but could not
replace source evidence to make a comparison pass.

Simulator credibility was assessed in Gate 3 by comparing eligible Python
trajectory endpoints with NASA's General Mission Analysis Tool (GMAT) under
the same accepted force-model configurations. An endpoint is the simulated
position and velocity at the end of a validation window. The diagnostic chain
isolated a formatting problem in the file containing the Earth-J2 coefficient
and confirmed the repaired format through degree-two, order-zero gravity
comparisons. In the final accepted evidence, all 67 evaluable checks passed;
the repaired endpoint set also passed the unchanged position and velocity
limits. This supports the implementation only within the tested public scope;
it does not validate a complete mission environment.

One record, the third return-trajectory-correction event (RTC3), is explicitly
excluded from the claim. Its qualified public trajectory file was created
before the event and therefore cannot show the event's reconstructed outcome
under the frozen temporal-evidence rule. It is
therefore classified as \texttt{not\_eligible}, rather than as a successful
or failed endpoint. This distinction prevents an important category error:
source eligibility is not a numerical outcome. The exclusion remains visible
in the evidence package rather than being silently omitted.

\subsection{Frozen development data and leakage control}

Let $\mathcal{D}=\{(x_i,y_i,z_i,g_i)\}_{i=1}^{N}$ denote the frozen
development dataset, where $N$ is the total number of development rows and $i$
indexes one simulated candidate plan. The vector $x_i$ contains only the
allowed scenario features. The target $y_i$ is robust total correction delta-v
in m/s. The binary label $z_i$ equals one when an independent propagation marks
the plan feasible and zero otherwise. Finally, $g_i$ identifies the related
mission-design group. Whole groups, rather than individual rows, are assigned
to one of five development folds. Consequently, related cases never appear on
both the training and validation sides of the same comparison.

All transformations are learned within the training partition of each fold.
These include filling missing values, converting categories to numbers,
rescaling features and targets, and principal-component analysis (PCA), a
linear compression of correlated features, where configured. Fitting any of
these steps to all rows would let information from the held-out validation
fold influence training. Preventing that leakage makes the comparison
interpretable.

The development campaign uses nested, deterministically selected sample rungs
of 128, 256, 512, and 1,024 training rows. Poorly performing configurations
stop at an early rung under a rule fixed before results; selected configurations
are then repeated with 20 frozen random seeds. These runs are not estimates of
an operational mission distribution. They are a prespecified stress test of
performance across group structure, sample size, and controlled randomness.
Calibration and final-test data were never read by these campaigns.

\subsection{Task definitions and estimands}

The primary endpoint is robust correction-cost prediction. For a fitted
predictor $f$, normalized root-mean-square error (NRMSE) is
\begin{equation}
\mathrm{NRMSE}(f) =
\frac{1}{s_y}\sqrt{\frac{1}{n}\sum_{i=1}^{n}\left(f(x_i)-y_i\right)^2},
\label{eq:nrmse}
\end{equation}
where $n$ is the number of held-out predictions, $f(x_i)$ is the predicted
delta-v for row $i$, $y_i$ is its simulated target, and $s_y$ is the sample
standard deviation of all frozen development targets. The numerator is the
usual root-mean-square error in m/s; division by $s_y$ makes NRMSE unitless and
comparable across folds. Zero is perfect and lower is better. For example,
NRMSE 0.01 means that the root-mean-square prediction error is 1\% of the
development target's standard deviation. The primary report pools predictions
made while each fold was held out (pooled out-of-fold NRMSE). Separate fold and
seed summaries are retained so that one favorable split or random run cannot
hide instability. Feasibility-constrained regret, the extra correction cost
caused by selecting a worse candidate plan, is reported as an interpretation
metric rather than as a replacement for the frozen primary endpoint.

The classification track estimates $p_f(x)\in[0,1]$, the predicted probability
that input $x$ is feasible under independent propagation. Its Brier score is
the mean squared probability error,
\begin{equation}
\mathrm{Brier}=\frac{1}{n}\sum_{i=1}^{n}\left(p_f(x_i)-z_i\right)^2,
\label{eq:brier}
\end{equation}
where $z_i\in\{0,1\}$ is the observed feasibility label; lower is better and
zero is perfect. Area under the receiver-operating-characteristic curve
(AUROC) measures ranking across all thresholds: 0.5 indicates chance-level
ranking and 1.0 indicates perfect ranking. At the fixed threshold $t=0.5$,
recall is
\begin{equation}
\mathrm{Recall}_{0.5}=\frac{\mathrm{TP}}{\mathrm{TP}+\mathrm{FN}}.
\label{eq:recall}
\end{equation}
Here TP counts actually feasible plans correctly marked feasible, and FN counts
actually feasible plans rejected by the classifier. Thus recall is the fraction
of feasible options retained. Precision is the fraction of plans marked
feasible that are actually feasible; an infeasible plan incorrectly admitted is
a false positive. This label direction matters: low recall reduces screening
utility, whereas unsafe acceptance must be assessed through precision,
false-positive behavior, and independent propagation. No one metric establishes
operational safety.

The Q01 family is described generically by an encoded-state kernel
\begin{equation}
K_q(x,x')=
\left|\left\langle 0^{\otimes q}\middle|
U_{\phi}(x)^{\dagger}U_{\phi}(x')
\middle|0^{\otimes q}\right\rangle\right|^2,
\label{eq:qkernel}
\end{equation}
where $x$ and $x'$ are two classical scenario vectors, $q\in\{4,6,8\}$ is the
number of simulated qubits, $|0^{\otimes q}\rangle$ is the all-zero initial
state, and $U_{\phi}(x)$ is the frozen circuit that encodes $x$. The symbol
$\dagger$ denotes the conjugate transpose, which reverses the circuit for the
overlap calculation. The squared overlap $K_q$ lies between zero and one: a
larger value means that two inputs look more similar under this particular
encoding. It is not a probability that either trajectory is feasible. The
pairwise values form a Gram matrix, a table of every input-to-input similarity,
which a kernel regressor converts into predictions. Exact features, PCA
dimensions, circuit settings $\phi$, bandwidths (how quickly similarity changes)
and regularization values (how strongly fitting is constrained) are preserved
in committed configuration records. Equation \eqref{eq:qkernel} defines a classically simulated research
representation, not a hardware implementation or quantum speedup.

\subsection{Models and matched controls}

The preregistered Q01 candidate is a fidelity-style quantum-kernel regressor
for correction cost. Its principal comparator is the classical physics-residual
model C06. If $b_{\mathrm{low}}(x)$ is a lower-fidelity physical estimate, C06
learns a correction $r(x)$ to form
\begin{equation}
\widehat{y}_{\mathrm{C06}}(x)=b_{\mathrm{low}}(x)+r(x).
\label{eq:c06}
\end{equation}
This is a stronger comparator than a generic regressor because it starts from
physical structure available in the benchmark and asks machine learning to
predict only the remaining error, or residual. The comparison set also contains
A01 random Fourier features, a finite classical approximation to a radial-basis
function (RBF) kernel; an RBF kernel treats nearby feature vectors as more
similar than distant ones. Compressed classical neural-network views and other
frozen configurations were also included. These controls test whether a
classical representation can explain an apparent quantum-feature effect
\cite{steffen2025,xie2022,rahimi2007}.

After the accepted Gate 5 result, exploratory protocol P001 defined two
follow-ups without reopening the primary campaign. Q01b replaces whole-state
overlap with projected features: it summarizes each simulated qubit through
local measurement expectations and applies a distance-based kernel to those
summaries. A deterministic Nystr\"om approximation limits the kernel to at most
256 representative training rows, called landmarks. FQK, whose repository ID
means \emph{feasibility-only quantum kernel}, applies the same projected-kernel
idea only to the independently propagated feasibility label; here FQK does not
mean ``fidelity quantum kernel.'' Both tracks retained groupwise splits,
training-fold-only preprocessing, nested rungs, matched controls, and 20 seeds.
They could produce new development-only evidence, but could not revise the
decision already made for Q01.

The 16 successor decisions D034--D049 tightened this discipline further. Each
protocol froze its candidate, target, comparator, threshold, stopping rule, and
source path before fitting. Their compact model names are immutable registry
identifiers; the repository stores each complete construction. Several
candidates produced a lower NRMSE than C06, but that alone was insufficient. A
candidate also had to beat a matched classical alternative given the same
inputs and stage structure. Such a control is often called a dequantization
control because it tests whether the apparent benefit survives when the
quantum-derived component is replaced by a classical one.

\subsection{Human-centered decision-support boundary}

The reported model outputs are treated as evidence for a future screening tool,
not as a command to change a trajectory. A human reviewer would need the model
estimate, the applicable scenario group, its comparator result, its uncertainty
record, and an explicit statement of whether the case is inside the authorized
scope. The research package preserves those items as source records and figure
provenance rather than reducing them to a single leaderboard value.

This stance also defines the fallback: when a candidate is ineligible, poorly
calibrated, source-ineligible, or outside the public benchmark boundary, the
model is not promoted and the physical or numerical reference remains the
decision basis. This is a design boundary for later human--AI interaction; it
does not evaluate an interface or the behavior of mission personnel.

\subsection{Statistical decision process}

The preregistered Gate 5 trigger required every eligible fold and seed to
complete, comparison with the strongest classical finalist and matched
dequantization controls, paired uncertainty checks, and at least one qualifying
regime. A \emph{regime} is a subgroup defined before outcomes, such as a
fidelity level or uncertainty family; it is not a favorable subgroup discovered
after inspection. A paired comparison subtracts model errors measured on the
same folds and seeds. The paired 95\% bootstrap interval repeatedly resamples
those matched differences to quantify uncertainty. A sign-permutation test
also asks whether the observed direction could arise if positive and negative
paired differences were exchangeable under no model advantage; multiplicity
control limits false discoveries across repeated comparisons. The report does
not add a post-outcome seed, model, split, or threshold to repair a weak result.
Learning curves and all retained seed outcomes are shown because training fit
or one favorable run is not adequate evidence of generalization
\cite{peters2023,grant2019}.

The distinction between task completion and eligibility is also retained. The
primary development campaign (decision D006) completed all 871 authorized
computational tasks with zero software or task failures. The variational
quantum regressor Q02 and hybrid quantum-residual model Q03 were not carried to
larger sample rungs because too few configurations met the frozen retention
rule at the first rung. Their status is therefore
\texttt{not\_reached\_under\_frozen\_eligibility}. It does not mean a
failed experiment, a missing computation, or a favorable result withheld from
the report.

\subsection{Computational-resource and reproducibility boundary}

The reference compute environment was a Windows laptop with a 13th Generation
Intel Core i9-13900HX central processing unit (CPU), 32~GB random-access memory
(RAM), and an NVIDIA GeForce RTX 4060 Laptop graphics processor (GPU). A
statevector stores all $2^q$ complex amplitudes of a $q$-qubit state in
classical memory. The experiment bounded this memory use and ran statevector
tasks one at a time; the GPU and quantum hardware were not used for the
reported QML evidence. The
hardware description is included so that readers understand the throughput and
resource constraints under which the public benchmark was executed. It is not
a prescription that future studies use identical hardware.

The source release contains configuration records, seeds, data cards, figure
registry entries, machine-readable tables, and reporting scripts. Every figure
used in this paper is copied with its corresponding registry entry into the
paper package. These measures make the evidence auditable while preserving the
lock on calibration and final-test data.
